% Generated draft: requires mathematical and source review.
\documentclass[11pt]{article}
\usepackage[margin=25mm]{geometry}
\usepackage{fontspec}
\setmainfont{Latin Modern Roman}
\usepackage{amsmath,amssymb,mathtools,unicode-math}
\setmathfont{Latin Modern Math}
\usepackage{enumitem,needspace,xurl,hyperref}
\setlength{\parindent}{0pt}
\setlength{\parskip}{5pt}
\setlength{\emergencystretch}{3em}
\newcommand{\N}{\mathbb{N}}
\newcommand{\Z}{\mathbb{Z}}
\newcommand{\Q}{\mathbb{Q}}
\newcommand{\R}{\mathbb{R}}
\newcommand{\C}{\mathbb{C}}
\newcommand{\F}{\mathbb{F}}
\newcommand{\A}{\mathbb{A}}
\newcommand{\PP}{\mathbb{P}}
\newcommand{\E}{\mathbb{E}}
\newcommand{\eps}{\varepsilon}
\newcommand{\dd}{\,\mathrm{d}}
\begin{document}
\section*{UnsolvedMath ID: 20002482}
\textbf{Source problem code:} catalogue entry
\subsection{Definitions and mathematical statement}
Paul Rabinowitz asks for a reasonable PDE analogue of Arnold diffusion and asks the same for infinite lattices such as the Fermi--Pasta--Ulam system. The source expresses the view that finite-dimensional transition chains may not be the right PDE mechanism and mentions Kuksin's high-Sobolev-norm ``diffusion'' results. It does not fix an equation, topology, orbit class, time scale, or definition of diffusion. The aid identifies the original number as Problem 25 rather than the extracted ``5''; the trailing $2$ is a page artifact.

The supplied literature account distinguishes arbitrary finite Sobolev amplification from unbounded growth along one trajectory. It reports finite amplification and several heteroclinic energy-exchange mechanisms for Hamiltonian PDEs, and genuine infinite transition chains for specially coupled infinite pendulum lattices. It reports unbounded Sobolev growth for the standard cubic defocusing NLS on the square torus, and diffusion for classical FPU, still open in the literature checked. These are supplied assessments without complete references or a check date; the special pendulum lattice is not silently identified with classical FPU.

The aid's proved calibration is kinematic. Let $a_k(t)$ be Fourier coefficients indexed by $k\in\Z^d$ with conserved positive mass
\[
M=\sum_k|a_k(t)|^2\in(0,\infty),\qquad
\mu_t=\frac1M\sum_k|a_k(t)|^2\delta_k.
\]
For $s\geq\tfrac12$ with finite relevant moments, define
\[
\|u(t)\|_{\dot H^s}^2=\sum_k|k|^{2s}|a_k(t)|^2,\qquad
R_s(t)=\left(\frac{\|u(t)\|_{\dot H^s}}{\sqrt M}\right)^{1/s}.
\]
For probability measures, $W_p(\mu,\nu)$ is the infimum over couplings of $(\int|k-\ell|^p\,\dd\pi(k,\ell))^{1/p}$, with $p=2s\geq1$. Then
\[
R_s(t)=W_{2s}(\mu_t,\delta_0),\qquad
|R_s(t)-R_s(0)|\leq W_{2s}(\mu_t,\mu_0)\leq R_s(t)+R_s(0).
\]
Thus for fixed initial data with finite $2s$-moment, unbounded homogeneous Sobolev growth is equivalent to unbounded $W_{2s}$ displacement of normalized Fourier mass. No PDE dynamics are assumed or constructed by these identities.

For $s>1$ and integers $N\to\infty$, the explicit thin-tail family is
\[
\mu_N=(1-N^{-s})\delta_0+N^{-s}\delta_{Ne_1}.
\]
It satisfies
\[
W_{2s}(\mu_N,\delta_0)=N^{1/2}\to\infty,
\quad W_1(\mu_N,\delta_0)=N^{1-s}\to0,
\quad\sup_A|\mu_N(A)-\delta_0(A)|=N^{-s}\to0.
\]
The last is the total-variation distance with the supremum-over-sets convention; an unhalved signed-measure norm is twice it. The moved fraction of mass vanishes even though the high moment diverges. Fourier amplitudes with these squared moduli have $\|u_N\|_{\dot H^s}=\sqrt M\,N^{s/2}$. This family is not asserted to be a trajectory of NLS, FPU, or any Hamiltonian PDE. It shows why high-norm growth alone does not force transport of a macroscopic mass fraction or a stronger dynamical version of Arnold diffusion.
\subsection{Short English statement}
The source asks broadly what Arnold diffusion should mean for PDEs and infinite lattices. The aid relates high Sobolev growth to a matching high-order Wasserstein displacement of modal mass, but a vanishing-mass tail shows that this need not produce macroscopic transport. The calibration does not construct a diffusing PDE orbit or settle classical FPU.
\subsection{Sources}
\begin{itemize}
\item Local UnsolvedMath registry, permanent ID $20002482$, \texttt{catalogue entry}; \texttt{aim-pdes-notes.json}, \texttt{notes[93]}, extracted source item 5, identified by the aid as original Problem 25.
\item American Institute of Mathematics, \emph{New connections between dynamical systems and PDE's}, Paul Rabinowitz contribution: \url{https://aimath.org/WWN/dynpde/dynpde.pdf}.
\item The source names Kuksin for high-Sobolev-norm diffusion. The aid additionally describes finite-amplification, heteroclinic and infinite-pendulum-chain results, without authors for those additional claims, titles, dates, or authoritative full citations.
\item Research attempt 1: partial result, problem partially solved at run, rigor label ``proved,'' novelty confidence low. Registry classification partial progress, status partially solved, statement status exact.
\end{itemize}
\end{document}
